% Integration-ready mathematical appendix. No standalone preamble.
\section{A different source: fixed coprime formula labels}
\label{sec:fermat-flow}

This section supplies a new source-level alternative to the affine-factor
and coefficient-budget sources inspected at commit
\texttt{c1a102936aca437e9999530a0111e6ddab56dc6d}. It does not supply
a smaller completed Diophantine equation. The change removes a freely
chosen common multiple or a coefficient-height budget from the source;
the cost of computing and retrieving its labels remains to be compiled.
The closure argument itself is the single-conclusion Horn-flow lemma
already used by package~68. The different ingredient here is a fixed,
parameter-free sequence of pairwise coprime labels.

\subsection{Coprime labels do not need to be primes}

For $j\in\mathbb N$ put
\[
 F_j=2^{2^j}+1.
\]
The elementary identity
\[
 \prod_{i<j}F_i=F_j-2
\]
follows by induction, including the empty product at $j=0$.
Consequently $F_j>1$ and the $F_j$ are pairwise coprime: a common divisor
of $F_i,F_j$, with $i<j$, divides $F_j$ and $F_j-2$, and is odd.
Primality is unnecessary; $F_5=641\cdot6700417$ is allowed.

For any finitely supported vectors $a,b\in\mathbb N^{(\mathbb N)}$,
\begin{equation}\label{eq:fermat-injective}
 \prod_jF_j^{a_j}=\prod_jF_j^{b_j}
 \quad\Longleftrightarrow\quad a=b.
\end{equation}
To prove the nontrivial implication, cancel the minimum of the two
exponents at every index. If a positive exponent remains on one side,
its factor $F_j>1$ is coprime to the whole other side, a contradiction.
This proof requires no factorization algorithm or primality predicate.

\subsection{An exact unordered proof certificate}

Let $\mathcal R$ be any family of finitary Horn rules with natural-number
labels, one conclusion $c_r$ and a finite premise multiset $P_r$.
Zero-premise rules are axioms. Let $D$ be the least rule-closed set.
Then for each target $t$ the following conditions are equivalent:
\begin{enumerate}
 \item $t\in D$.
 \item Some finite list of valid rule instances satisfies
 \begin{equation}\label{eq:fermat-product}
  \prod_{i<N}F_{c_i}
    =F_t\prod_{i<N}\prod_{p\in P_i}F_p.
 \end{equation}
 \item Some finite list of valid rule instances and natural numbers
 $S_0,\ldots,S_N$ satisfy
 \begin{equation}\label{eq:fermat-accumulator}
  S_0=F_t,\qquad S_N=1,\qquad
  S_{i+1}F_{c_i}=S_i\prod_{p\in P_i}F_p.
 \end{equation}
\end{enumerate}

For completeness, expand a finite proof tree root first. Let $S_i$ be
the product of the weights of the current open obligations. Expanding a
node replaces its conclusion by its actual premises, giving
\eqref{eq:fermat-accumulator}; the empty frontier has product one.

For soundness of the product certificate,
\eqref{eq:fermat-injective} gives the coordinate balance
\[
 \sum_{i<N}\left(e_{c_i}-\sum_{p\in P_i}e_p\right)=e_t.
\]
If $t\notin D$, sum the coordinates outside $D$. Each rule contributes
at most zero: if its conclusion is outside $D$, at least one premise is
outside $D$, while a conclusion inside $D$ contributes zero minus a
nonnegative number. The right side is one, a contradiction. Thus cycles
cannot create a false proof. A disconnected cycle can be present in a
valid certificate, but it is harmless. Saturating the finitely many
selected rules gives a finite proof DAG for the target; this can be
unfolded to a proof tree.

The accumulator conditions imply the product certificate by multiplying
and cancelling. All $S_i$ are positive: this follows inductively from
$S_0>0$ and positivity of every factor in
\eqref{eq:fermat-accumulator}. No independent positivity guard is needed
when the accumulators range over naturals. We do not claim that an
arbitrary ordering of an arbitrary balanced multiset gives an integral
accumulator path. Existence is enough: the recovered proof tree provides
another valid list and path if necessary.

Exactly one conclusion and nonnegative rule multiplicities are essential.
The rule $a\longrightarrow\{a,t\}$ has balance $e_t$ without any axiom.
Also, subtracting the vector for $a,a\longrightarrow b$ from that for
$a\longrightarrow b$ gives $e_a$, although neither rule derives anything
from the empty set. Those are excluded by the stated hypotheses.

For the repository's target-one calculus, use the actual arity of each
rule, so that truth code zero is a genuine zero-premise axiom. Then
$F_0=3$ and the initial target weight is the fixed integer $F_1=5$.
One must not simultaneously use $F_0=3$ and treat padded truth premises
as free factors. A different, also sound convention sets the zero-label
weight to one and retains Fermat weights only for positive labels, but
requires a separate zero-label branch. The convention here avoids it.

\subsection{A precise four-parameter bounded source}

The double exponent in $F_j$ is not an ordinary polynomial. To make its
remaining compilation cost concrete, define
\[
 \beta(x,y,j)=x\bmod(1+(j+1)y).
\]
Use four outer natural parameters $N,x_S,x_Q,y$, and put
$S_i=\beta(x_S,y,i)$, $Q_i=\beta(x_Q,y,i)$.
Impose the endpoints
\[
 S_0=5,\qquad S_N=1,\qquad Q_0=2.
\]
For every $i<N$, require simultaneously
\[
 Q_{i+1}=Q_i^2
\]
and one actual rule instance $P_i\longrightarrow c_i$, all of whose
labels are strictly less than $N$, with
\[
 S_{i+1}\bigl(\beta(x_Q,y,c_i)+1\bigr)
   =S_i\prod_{p\in P_i}\bigl(\beta(x_Q,y,p)+1\bigr).
\]
Allow identity rules $c\longrightarrow c$ as padding; their addition
does not change derivability. These conditions use one bounded universal
block, with existential local witnesses. Induction gives
$Q_j=2^{2^j}$ for $0\le j\le N$, so every referenced label is authenticated.
The preceding theorem proves soundness.

For completeness, choose a finite proof tree and then a horizon $N$
larger than all labels and at least its number of nodes. Pad the rule list
by identities, keeping the terminal accumulator equal to one. Form the
finite $S$ and $Q$ sequences first. Choose
\[
 y=N!\left(1+\max_{0\le i\le N}\{S_i,Q_i\}\right).
\]
The moduli $1+(i+1)y$ are pairwise coprime and exceed the desired entries.
The finite Chinese remainder theorem supplies $x_S,x_Q$ with the same
scale $y$. This factorial is used only to prove existence of CRT witnesses;
it is not a function in the proposed final polynomial, and no source
condition asserts that an unknown is a factorial. The choice order avoids
any self-referential size bound.

\subsection{What must still be charged}

The preceding four-parameter source is a finite first-order arithmetic
formula, not an ordinary fixed-dimensional polynomial. Its costs include
the supplied Horn-rule recognizer, the two adjacent sequence entries,
up to three indexed label lookups per binary rule, label-range witnesses,
and the bounded-universal compiler. For example a lookup
$u=\beta(x_Q,y,p)$ has the completely explicit quadratic graph
\[
 d=(p+1)y,\qquad x_Q=(d+1)k+u,\qquad u+s=d
\]
over naturals. These are three equations using four local witnesses
$d,k,u,s$, one of which is the returned value $u$. A strict label bound
is $p+b+1=N$ with one further natural witness $b$.
Shared denominators and fixed zero-premise branches may save registers,
but no unmeasured savings are assumed.

Alternatively, one can define a label by two separately compiled
exponentiation graphs. That removes the power sequence and introduces the
corresponding graph witnesses. Both routes remove the source's common
multiple or coefficient budget; neither has yet been shown to reduce the
complete polynomial's variable count, degree or arithmetic cost.
The finite script \texttt{alternative/verify.py} checks all rule-multiset cases
listed in \texttt{alternative/VERIFICATION.json}, constructs proof trees and
integer traces, verifies a small shared-CRT example, and includes the
failure modes above. It is an adversarial regression, not a substitute
for these universal proofs or the inherited calculus-to-ZFC calibration.
